\section{Introduction}
\label{sec:intro}
% PLAN (about 1.5 pages, five paragraphs)
%   P1 Clinical problem. Ribs and clavicles overlap lung lesions on chest radiographs; dual-energy subtraction (DES) is the reference but needs
%      special hardware and extra exposure (sun2026bsldm, sun2026debonedit). Software suppression is attractive; reader and CAD studies report
%      benefit for nodule detection (freedman2011lung, schalekamp2014cad).
%   P2 CT-supervised suppression. Digitally reconstructed radiographs from chest CT give paired supervision without DES (angaitkar2026anatomy);
%      diffusion models trained on real DES pairs are the other line (chen2024bsdiff, sun2026bsldm, sun2025gllcm, sun2026debonedit).
%   P3 The gap. Suppressors are judged by fidelity to a bone or soft-tissue reference (PSNR, SSIM). A model that is rewarded for predicting bone
%      can also absorb a lesion into the bone image. A suppressor that does nothing never absorbs a lesion, so absorption must be read together
%      with suppression gain. No counterfactual test exists in the bone-suppression literature we know of [novelty claim only after the
%      systematic search listed in paper/REFS_TODO.md].
%   P4 Contributions (C1-C4): (i) counterfactual lesion audit with an exact-truth version on synthetic data and a confirmatory version on real
%      radiographs with independent rib masks; (ii) twin lesion-invariance regulariser, with a dose-response over its weight and an ablation of
%      its parts; (iii) a public-data recipe with a K=3 ensemble, uncertainty and a modern-encoder screen; (iv) a locked-set benchmark with
%      downstream detection and classification, a pre-specified plan and a licensed release.
%   P5 Outline of the paper.
Ribs and clavicles overlap the lung fields on a chest radiograph and can obscure abnormalities \citep{angaitkar2026anatomy}.
Dual-energy subtraction (DES) separates bone from soft tissue, but it needs dedicated hardware and a second exposure
\citep{kuhlman2006des,sun2026bsldm,sun2026debonedit}. Software bone suppression avoids both, and reader and CAD studies have examined its effect on nodule
detection \citep{freedman2011lung,schalekamp2014cad}.

Two lines of learning-based suppression are current. One renders paired supervision from chest CT: each scan is split into bone, non-lung
soft tissue and lung, the parts are projected separately, and a network predicts one component and recovers the other by subtraction
\citep{angaitkar2026anatomy}. The other trains diffusion models on real DES pairs
\citep{chen2024bsdiff,sun2026bsldm,sun2025gllcm,sun2026debonedit}. In the CT-supervised line, lesion preservation has been assessed indirectly,
through detection gains and the contrast of boxed lesions \citep{angaitkar2026anatomy}; we measure it directly, with a known inserted signal, and
build the measurement into training.

Suppressors are judged by how closely the output matches a bone or soft-tissue reference. That judgement misses a failure that matters
clinically: a network rewarded for predicting bone can also move lesion signal into the bone image, and the lesion then disappears from the
suppressed radiograph. Because a model that does nothing never absorbs a lesion, absorption has to be read next to suppression gain.
Inserted or simulated nodules have been used before to measure contrast after suppression, with a hand-built clavicle suppression method and
with phantom images \citep{hogeweg2013suppression,takagi2020dose}, and the loss or hallucination of image content by learned image-to-image
models is documented in other modalities \citep{cohen2018distribution,antun2020instabilities,bhadra2021hallucinations,tivnan2024hallucination}.
In our search (PubMed, arXiv and Crossref, abstracts read) we found no audit of learned bone suppression against an exactly known, additive
lesion signal, and no training objective aimed at it.

We make four contributions.
\begin{enumerate}
\item A counterfactual lesion audit that measures the fraction of an inserted lesion a suppressor moves into its bone image, reported beside
 suppression strength, with sites shared by all models and bone overlap taken from independent annotations.
\item A twin lesion-invariance regulariser that ties the bone predictions of an image and of a copy with a pasted lesion, with a dose-response over
 its weight and an ablation of its parts.
\item A training recipe on public CT-derived projections only, with a three-model ensemble, split-conformal intervals and a screen of modern encoders.
\item A pre-specified evaluation on locked sets (paired dual-energy images, an annotated rib set, real nodules and four downstream tasks) and an
 external validation on paediatric radiographs, with every change to the analysis disclosed.
\end{enumerate}

Section~\ref{sec:related} reviews related work. Section~\ref{sec:data} describes the data, Sections~\ref{sec:methods} and~\ref{sec:method_v2} the
methods, Section~\ref{sec:audit} the audit, Section~\ref{sec:experiments} the experiments and Section~\ref{sec:results} the results; we then discuss
limitations and reproducibility.
